\documentclass[11pt]{article}

\usepackage[margin=1in]{geometry}
\usepackage[T1]{fontenc}
\usepackage{lmodern}
\usepackage{microtype}
\usepackage{amsmath,amssymb,amsthm,bm}
\usepackage{booktabs,tabularx,multirow}
\usepackage{enumitem}
\usepackage[authoryear,round]{natbib}
\usepackage{xcolor}
\usepackage[colorlinks=true,linkcolor=blue!50!black,citecolor=blue!50!black,urlcolor=blue!50!black]{hyperref}

\newtheorem{assumption}{Assumption}
\renewcommand{\theassumption}{H\arabic{assumption}}
\newtheorem{proposition}{Proposition}
\newtheorem{lemma}{Lemma}
\theoremstyle{definition}
\newtheorem{definition}{Definition}
\theoremstyle{remark}
\newtheorem{remark}{Remark}

\newcommand{\E}{\mathbb{E}}
\newcommand{\Var}{\operatorname{Var}}
\newcommand{\Cov}{\operatorname{Cov}}
\newcommand{\Corr}{\operatorname{Corr}}
\newcommand{\cum}{\operatorname{cum}}
\newcommand{\HSIC}{\operatorname{HSIC}}
\newcommand{\indep}{\perp\!\!\!\perp}
\newcommand{\xexo}{x_{\mathrm{exo}}}
\newcommand{\tx}{\tilde{x}}
\newcommand{\ty}{\tilde{y}}

\title{Synthetic Instruments by Adversarial Splitting:\\
Identification through Non-Gaussianity\thanks{Research proposal and Monte Carlo testbed.
Code, tests and simulation scripts accompany this paper in the \texttt{ivsplit} package
(\texttt{src/ivsplit}, \texttt{tests}, \texttt{scripts}). The literature search is not exhaustive.}}
\author{Draft}
\date{October 19, 2026}

\begin{document}
\maketitle

\begin{abstract}
We study the estimation of a linear structural coefficient when a regressor is endogenous and no
credible external instrument is available. The idea is to split the endogenous regressor into an
exogenous component and a remainder, $x = \xexo + r$ with $\xexo \indep u$, and to use $\xexo$ as an
internally generated (``synthetic'') instrument or, equivalently, $r$ as a control function. We
first explain why a naive version of this idea---generating many candidate instruments with a
generative network trained to pass weak-instrument and over-identification tests---cannot work:
exogeneity never enters a computable criterion. We then make the identifying assumption explicit:
the regressor and the structural error are driven by mutually independent sources, at least two of
which are non-Gaussian. Under this assumption the structural coefficient is identified from
fourth-order cumulants up to a two-point set $\{\beta,\beta+\lambda\}$, a label ambiguity that is
intrinsic to the model and that we resolve with a sign restriction or an auxiliary variable. The
estimation layer is a neural network that learns the split under a Hilbert--Schmidt independence
penalty, trained on one half of the sample, with estimation and weak-identification-robust
inference on the other half. We also show that, in the over-complete linear model, the exogenous
source is generally \emph{not} a measurable function of the data, so the network can only target
moment validity rather than exact independence; a pilot simulation confirms that the unanchored
split collapses towards OLS while the proposed diagnostics correctly report the failure. We
propose a complete battery of pre-estimation, relevance, exogeneity and sensitivity tests, and a
Monte Carlo design that pits the method against OLS, an oracle instrument, structural error
projection, heteroskedasticity-based instruments, Gaussian-copula control functions and a linear
cumulant (over-complete ICA) estimator.
\end{abstract}

\noindent\textbf{Keywords:} endogeneity, internal instruments, independent component analysis,
higher-order cumulants, HSIC, control function, weak identification.

\section{Introduction}\label{sec:intro}

Consider the linear structural model
\begin{equation}\label{eq:model}
  y = x\beta + w'\delta + u, \qquad \E[xu]\neq 0,
\end{equation}
where $x$ is a scalar endogenous regressor, $w$ is a vector of exogenous covariates and $u$ is the
structural error. The textbook solution is an external instrument: a variable that moves $x$ and is
excluded from \eqref{eq:model}. When no such variable is credible, researchers turn to
\emph{internal} instruments, built from the data at hand by exploiting heteroskedasticity
\citep{rigobon2003,lewbel2012}, higher moments \citep{lewbel1997,erickson2002}, latent discrete
structure \citep{ebbes2005}, distributional assumptions \citep{parkgupta2012} or functional form
\citep{dong2026sp}.

This paper develops a variant of that programme. We seek a decomposition
\begin{equation}\label{eq:split}
  x = \xexo + r, \qquad \xexo \indep u,
\end{equation}
and use $\xexo$ as the source of exogenous variation. Two ingredients define the proposal:
\begin{enumerate}[label=(\roman*)]
  \item \textbf{Identification} comes from mutually independent, non-Gaussian sources of variation,
  as in independent component analysis \citep[ICA;][]{comon1994} and its causal counterparts with
  latent variables \citep{hoyer2008,salehkaleybar2020,wang2023,tramontano2025}.
  \item \textbf{Estimation} uses a neural network that learns the split, with independence enforced
  by a Hilbert--Schmidt independence criterion \citep[HSIC;][]{gretton2005,gretton2008} in the
  spirit of HSIC-X \citep{saengkyongam2022}, and with sample splitting so that inference does not
  reuse the data that shaped the instrument.
\end{enumerate}

Our contributions are fourfold. First, we state precisely what has to be assumed: any synthetic
instrument replaces the exclusion restriction of an external instrument with another untestable
hypothesis, here independence of latent sources (Section~\ref{sec:why}). Second, we give a
self-contained identification result for the canonical latent-confounder model, which shows that
the coefficient is identified from fourth-order cumulants only up to the set $\{\beta,\beta+\lambda\}$
and that this ambiguity is a genuine observational equivalence, not an artefact of the estimator
(Section~\ref{sec:identification}). Third, we show that the exogenous source is generally not a
measurable function of observables in this over-complete setting, which sharpens what a network can
and cannot learn and motivates an anchored version of the estimator (Section~\ref{sec:network}).
Fourth, we provide a complete testing battery and a Monte Carlo design, with an open-source
implementation and unit tests (Sections~\ref{sec:tests}--\ref{sec:mc}).

\section{Why the naive synthetic-instrument idea fails}\label{sec:why}

The first version of this project generated $2k$ candidate instruments with a generative
adversarial network trained to pass weak-instrument and Sargan tests, and filtered the excess with
Hausman tests. The design fails for reasons that are worth recording because they constrain every
variant of the idea.

\begin{enumerate}[label=\arabic*.]
  \item \emph{Exogeneity is not in the objective.} The condition $\Cov(z,u)=0$ involves the
  unobserved $u$ and appears in no computable criterion. A search over candidates converges to
  \emph{strong} instruments, not to \emph{valid} ones.
  \item \emph{Passing an over-identification test is necessary, not sufficient.} Instruments
  manufactured from the same data inherit the same contamination and agree with each other on the
  same wrong value; the Sargan--Hansen test \citep{sargan1958,hansen1982} has no power against a
  common bias.
  \item \emph{Training to pass a test is selection on the test outcome.} It produces pre-test bias,
  a winner's curse and a severe multiple-testing problem across $2k$ candidates.
  \item \emph{The Hausman test answers a different question.} \citet{hausman1978} compares OLS with
  IV and tests the endogeneity of the regressor. Testing a subset of instruments calls for the
  difference-in-Hansen (C) statistic \citep{eichenbaum1988}.
  \item \emph{No algorithm creates identification.} If the data contain no source of exogenous
  variation, no transformation of them will. Every synthetic instrument trades the exclusion
  restriction of an external instrument for another untestable hypothesis.
\end{enumerate}

The question therefore becomes: \emph{which identifying hypothesis is the researcher willing to
maintain in order to perform the split \eqref{eq:split}?} We answer it with independence and
non-Gaussianity of latent sources.

\section{Model and identification}\label{sec:identification}

\subsection{The canonical latent-confounder model}

Let $\tx = x - \E[x\mid w]$ and $\ty = y - \E[y\mid w]$ denote the regressor and outcome after
partialling out the exogenous covariates. Our leading structure is
\begin{equation}\label{eq:canonical}
  x = g(w) + H + e_X, \qquad u = \lambda H + e_Y, \qquad y = \beta x + w'\delta + u,
\end{equation}
so that $\tx = H + e_X$ and $\ty = \beta\tx + u = (\beta+\lambda)H + \beta e_X + e_Y$. Here $H$ is a
latent confounder, $e_X$ is the exogenous variation in $x$ that we wish to recover, and $e_Y$ is
idiosyncratic outcome noise. All three are mean zero and independent of $w$.

\subsection{Assumptions}

\begin{assumption}[Independent sources]\label{ass:indep}
$H$, $e_X$ and $e_Y$ are mutually independent and independent of $w$.
\end{assumption}

\begin{assumption}[Non-Gaussianity]\label{ass:nongauss}
The fourth cumulants $\kappa_H = \cum_4(H)$ and $\kappa_X = \cum_4(e_X)$ are non-zero.
\end{assumption}

\begin{assumption}[Class restriction]\label{ass:class}
The map from sources to observables belongs to a class for which the independent decomposition is
unique up to the symmetries made explicit below: linear mixing after partialling out $w$ as in
\eqref{eq:canonical}, or a non-linear class made identifiable by additive noise, sparsity of the
mixing map or an observed auxiliary variable
\citep{hyvarinen2019,khemakhem2020,zheng2022}.
\end{assumption}

\begin{assumption}[Relevance]\label{ass:relevance}
The exogenous component has non-zero variance and is correlated with $x$: $\Var(e_X)>0$.
\end{assumption}

\begin{assumption}[Labelling]\label{ass:label}
One of the following holds.
\textbf{(L1)} The sign of the confounding, $\operatorname{sign}(\lambda)$, is known.
\textbf{(L2)} There is an auxiliary variable $a$ (a component of $w$ or an additional observable)
such that the law of $H$ does not depend on $a$ while $\Var(e_X\mid a)$ does.
\end{assumption}

Assumption~\ref{ass:indep} is a disguised exclusion restriction and cannot be tested.
Assumption~\ref{ass:nongauss} is checkable. Assumption~\ref{ass:class} is necessary because
non-linear ICA is not identifiable in general \citep{hyvarinen1999}. Assumption~\ref{ass:label} is
needed because of the symmetry established next.

\subsection{Identification up to a two-point set}

Write $c_{4-k,k} = \cum(\underbrace{\tx,\dots,\tx}_{4-k},\underbrace{\ty,\dots,\ty}_{k})$ for the joint
fourth-order cumulants, and abbreviate $c_k \equiv c_{4-k,k}$. By multilinearity of cumulants and
independence of the sources,
\begin{equation}\label{eq:cumulants}
  c_k = \kappa_H (\beta+\lambda)^k + \kappa_X \beta^k + \kappa_Y\,\mathbf{1}\{k=4\}, \qquad k=0,\dots,4.
\end{equation}
For $k\le 3$ this is a two-term exponential sum in the unknown slopes $s_H=\beta+\lambda$ and
$s_X=\beta$, so it satisfies the linear recursion of Prony's method:
\begin{equation}\label{eq:prony}
  c_{k+2} = p\,c_{k+1} + q\,c_k,\quad k=0,1, \qquad p = s_H + s_X,\; q = -s_H s_X.
\end{equation}

\begin{proposition}[Identification up to labels]\label{prop:twopoint}
Under Assumptions~\ref{ass:indep}--\ref{ass:nongauss} in model \eqref{eq:canonical} with
$\lambda\neq0$:
\begin{enumerate}[label=(\alph*)]
  \item The matrix $M = \begin{pmatrix} c_{1} & c_{0}\\ c_{2} & c_{1}\end{pmatrix}$ has
  $\det M = -\kappa_H\kappa_X\lambda^2 \neq 0$, so $(p,q)$ is identified from
  $M(p,q)' = (c_2,c_3)'$, and the set $\{\beta,\beta+\lambda\}$ is identified as the roots of
  $t^2 - pt - q = 0$.
  \item The set is sharp: the model with parameters $(\beta',\lambda') = (\beta+\lambda,-\lambda)$ and
  sources $(H',e_X') = (e_X,H)$ generates exactly the same distribution of $(\tx,\ty)$.
\end{enumerate}
\end{proposition}

\begin{proof}
(a) Substituting \eqref{eq:cumulants},
$\det M = (\kappa_H s_H + \kappa_X s_X)^2 - (\kappa_H s_H^2 + \kappa_X s_X^2)(\kappa_H + \kappa_X)
= -\kappa_H\kappa_X (s_H - s_X)^2 = -\kappa_H\kappa_X\lambda^2$. The recursion \eqref{eq:prony}
follows from $s_j^{k+2} = p s_j^{k+1} + q s_j^{k}$ for $j\in\{H,X\}$, which holds because each $s_j$
solves $t^2 - pt - q = 0$. (b) Under the alternative parameters,
$\tx = H' + e_X' = e_X + H$ and $\ty = (\beta'+\lambda')H' + \beta' e_X' + e_Y
= \beta e_X + (\beta+\lambda)H + e_Y$, which coincide with the original variables.
\end{proof}

\begin{remark}
Part (b) is the population version of the ``identified only up to two points'' result of
\citet{delbianco2026} and of the label ambiguity in latent-variable LiNGAM
\citep{hoyer2008,tramontano2025}. It is not a weakness of any particular estimator: without
Assumption~\ref{ass:label} \emph{no} procedure can tell the confounder from the exogenous source,
because both enter $x$ with unit loading.
\end{remark}

\begin{proposition}[Labelling]\label{prop:label}
Under the conditions of Proposition~\ref{prop:twopoint}: (L1) if $\lambda>0$ then $\beta$ is the
smaller root, and if $\lambda<0$ the larger; (L2) $\beta$ equals the slope of
$\Cov(\tx,\ty\mid a)$ on $\Var(\tx\mid a)$ across values of $a$.
\end{proposition}

\begin{proof}
(L1) is immediate. For (L2), $\Var(\tx\mid a) = \sigma_H^2 + \sigma_X^2(a)$ and
$\Cov(\tx,\ty\mid a) = (\beta+\lambda)\sigma_H^2 + \beta\sigma_X^2(a)$, so the two conditional moments
are affine in $\sigma_X^2(a)$ with slope ratio $\beta$.
\end{proof}

\begin{remark}
Under (L2) the second-moment argument alone identifies $\beta$, without any appeal to
non-Gaussianity; it is identification through heteroskedasticity
\citep{rigobon2003,lewbel2012}. In our framework (L2) is used only to \emph{pick} one of the two
roots, which needs far less than consistent estimation of the slope; the cumulant equations then
provide an over-identifying check.
\end{remark}

\begin{proposition}[Gaussian non-identification]\label{prop:gauss}
If $H$, $e_X$ and $e_Y$ are Gaussian, $\beta$ is not identified: the distribution of $(\tx,\ty)$
depends on the five parameters $(\beta,\lambda,\sigma_H^2,\sigma_X^2,\sigma_Y^2)$ only through the
three moments $\Var(\tx)$, $\Cov(\tx,\ty)$ and $\Var(\ty)$, and every $\beta'$ in a non-degenerate
interval containing $\beta$ is observationally equivalent.
\end{proposition}

\begin{proof}
Gaussian vectors are characterised by their covariance. Fix $\beta'$ near $\beta$ and keep
$\sigma_H'^2 = \sigma_H^2$, $\sigma_X'^2=\sigma_X^2$. Choose $\lambda'$ so that
$\Cov(\tx,\ty) = \beta'\Var(\tx) + \lambda'\sigma_H^2$ and $\sigma_Y'^2$ so that $\Var(\ty)$ is
matched; by continuity $\sigma_Y'^2>0$ for $\beta'$ close enough to $\beta$.
\end{proof}

Proposition~\ref{prop:gauss} is the population reason why scenario~3 of the Monte Carlo is a
negative control: an honest procedure must report weak or failed identification there, rather than
a confident estimate. It also explains why the identification-strength statistic of
Section~\ref{sec:diagnostic} is built on $\det M$.

\subsection{An identification-strength diagnostic}\label{sec:diagnostic}

Proposition~\ref{prop:twopoint} shows that identification hinges on $\det M = -\kappa_H\kappa_X
\lambda^2$, which vanishes when either source is mesokurtic (e.g.\ Gaussian) or when there is no
confounding. We report
\begin{equation}\label{eq:D}
  D = \frac{\widehat{\det M}^{\,2}}{\widehat{\Var}_{\mathrm{boot}}\big(\widehat{\det M}\big)},
\end{equation}
a Wald statistic for $H_0:\det M = 0$, with a nonparametric bootstrap variance. Small values of $D$
(below the $\chi^2_1$ critical value) indicate that the higher-order information is insufficient to
separate the sources, in the spirit of the $D$ diagnostic of \citet{delbianco2026}. When $\lambda=0$
a small $D$ is benign---OLS is then consistent---and should be read together with a Hausman test.

\subsection{The linear identification layer as an estimator}

Replacing population cumulants by sample cumulants of the $w$-residualised data, solving the
(possibly over-identified, if third-order cumulants are added) linear system for $(\hat p,\hat q)$ and
labelling the roots with (L1) or (L2) gives a simple closed-form estimator, which we call
\emph{Cumulant-ICA}. It is the over-complete ($2\times3$) ICA estimator of the canonical model and
is closely related to the higher-moment estimators of \citet{lewbel1997} and \citet{erickson2002}.
We use it both as a competitor and as an anchor for the network (Section~\ref{sec:anchored}).

\section{The neural split}\label{sec:network}

\subsection{Procedure}

\begin{enumerate}[label=\textbf{Step \arabic*.},leftmargin=*]
  \item \textbf{Sample splitting.} Randomly split the sample into halves $A$ and $B$. All learning
  happens on $A$; $\hat\beta$, standard errors and every test are computed on $B$.
  \item \textbf{Partialling out.} Compute $\tx,\ty$ by regressing $x,y$ on $w$ (linearly, with a
  sieve, or with cross-fitted machine learning as in \citealp{chernozhukov2018}).
  \item \textbf{Training.} A network $f_\theta$ maps $(\tx, w)$ to $\xexo$ and sets
  $r = \tx - \xexo$. The parameters minimise
  \begin{equation}\label{eq:loss}
  \mathcal{L}(\theta,\beta) = \alpha\,\HSIC(\xexo, r)
     + \gamma\,\HSIC(\ty - \beta\tx,\, \xexo)
     + \eta\,\big[(c_{\min} - \rho^2)_+ + (\rho^2 - c_{\max})_+\big],
  \end{equation}
  where $\rho = \Corr(\xexo,\tx)$. The first term imposes independence between the two parts of the
  split; the second is the HSIC-X criterion of \citet{saengkyongam2022} with the learned
  instrument; the third is a relevance band that rules out degenerate solutions. The
  reconstruction term $\|x - \xexo - r\|^2$ of the original proposal vanishes by construction.
  $\beta$ is updated in alternation by the IV formula
  $\beta \leftarrow \sum \xexo\ty / \sum\xexo\tx$.
  \item \textbf{Estimation.} On $B$, compute $\xexo$ and $r$ and estimate $\beta$ by 2SLS with
  $\xexo$ as the excluded instrument (``split-IV''), or by OLS of $y$ on $(x,w,r)$ (``split-CF'').
  \item \textbf{Inference.} Report heteroskedasticity-robust standard errors, the robust
  first-stage $F$ \citep{montielolea2013}, and the Anderson--Rubin confidence set
  \citep{andersonrubin1949}, which remains valid when the learned instrument is weak.
  \item \textbf{Diagnostics.} Report $D$ from \eqref{eq:D} on $B$ next to $\hat\beta$, and the
  battery of Section~\ref{sec:tests}.
\end{enumerate}

\subsection{Degenerate solutions}

Three degenerate solutions must be excluded. (i) $\xexo\approx$ constant is independent of
everything but irrelevant; the relevance floor $c_{\min}$ excludes it. (ii) $\xexo = \tx$ has
maximal relevance and, with $r\equiv 0$, trivially satisfies $\HSIC(\xexo,r)=0$, but the IV estimator
collapses to OLS; the relevance cap $c_{\max}<1$ excludes it. (iii) Over-fitting on the training
sample inflates the first-stage $F$ and recreates the many-instruments bias; sample splitting
neutralises it.

\subsection{What the network can and cannot recover}\label{sec:cannot}

A natural reading of \eqref{eq:split} is that the network should output $\xexo = e_X$. The next
result shows that this is generally impossible in the over-complete model, where three sources
drive two observables.

\begin{proposition}[No measurable exact split]\label{prop:nosplit}
In model \eqref{eq:canonical} under Assumption~\ref{ass:indep}, suppose $\lambda\neq0$, the
characteristic function of $e_Y$ has no real zeros, $H$ has support $\mathbb{R}$, and the
characteristic function of $e_X$ has no real zeros. Then every measurable $f$ with
$f(\tx)\indep u$ is almost surely constant.
\end{proposition}

\begin{proof}[Proof sketch]
Since $e_Y$ is independent of $(H,e_X)$ and its characteristic function does not vanish,
$f(\tx)\indep \lambda H + e_Y$ implies $f(\tx)\indep H$ by deconvolution of the joint characteristic
function. For any Borel set $A$, independence then requires
$\Pr(H+e_X\in A\mid H=h)=\int \mathbf{1}_A(h+t)\,dF_{e_X}(t)$ to be constant in $h$ on the support
of $H$, i.e.\ the convolution of $\mathbf{1}_A$ with the law of $-e_X$ is constant on $\mathbb{R}$.
Taking Fourier transforms and using that the characteristic function of $e_X$ has no zeros, by
Wiener's Tauberian theorem $\mathbf{1}_A$ is a.e.\ constant, so the $\sigma$-field generated by
$f(\tx)$ is trivial.
\end{proof}

Laplace, Gaussian and centred exponential laws all satisfy the conditions, and adding $w$ to the
inputs does not help when $w$ is independent of the sources. Two consequences follow.

\begin{itemize}
  \item The population objective \eqref{eq:loss} cannot be driven to zero, and its minimiser in
  $(\theta,\beta)$ need not be at the true $\beta$. In particular, at the OLS slope the residual
  $\ty - \beta_{\mathrm{OLS}}\tx$ is uncorrelated with $\tx$, so the HSIC-X term only penalises
  higher-order dependence, which is weak in moderate samples. Our pilot simulations
  (Section~\ref{sec:pilot}) confirm that the unanchored split converges near OLS.
  \item What \emph{is} attainable is \emph{moment} validity. Writing $m(\tx)=\E[H\mid\tx]$, any
  function $f$ with $\E[f(\tx)m(\tx)]=0$ and $\E[f(\tx)\tx]\ne0$ is a valid and relevant instrument,
  since $\E[f(\tx)u] = \lambda\E[f(\tx)m(\tx)]$. Such functions exist in abundance (they form a
  subspace of codimension one), but selecting one requires knowledge of $m$, which in turn depends on
  the laws of $H$ and $e_X$---exactly what the cumulant layer identifies up to labels.
\end{itemize}

There are also non-generic cases in which an exact split exists, for example when $H$ is discrete
and $e_X$ has bounded support with non-overlapping translates; such structures are the province of
latent-mixture approaches \citep{ebbes2005}.

\subsection{Anchored split}\label{sec:anchored}

Proposition~\ref{prop:nosplit} motivates a revised estimation layer that lets the identification
layer do the identifying and the network do the approximating:
\begin{enumerate}[label=(\roman*)]
  \item On sample $A$, estimate $\beta$ (and, if desired, the source laws) with the cumulant
  layer and label the roots with (L1) or (L2), obtaining $\hat\beta_A$.
  \item Train $f_\theta$ on $A$ with $\beta$ fixed at $\hat\beta_A$ in the HSIC-X term, so that the
  network searches for instrument functions $f(\tx,w)$ whose implied moment conditions are
  approximately satisfied at the identified value, and use the relevance band to keep them
  informative.
  \item On $B$, estimate $\beta$ by 2SLS with the learned functions and test the over-identifying
  restrictions they imply.
\end{enumerate}
The non-linear value of the network lies where the linear layer is silent: a non-linear
first stage $g(w)$, auxiliary variables that modulate the law of $e_X$ (Assumption~\ref{ass:class}),
and efficient combination of moment conditions. Whether the anchored split improves on the plain
cumulant estimator in finite samples is an empirical question that the Monte Carlo is designed to
answer.

\section{Testing battery}\label{sec:tests}

All tests are computed on the estimation half $B$. Table~\ref{tab:tests} summarises them.

\begin{table}[ht]
\centering
\small
\caption{Tests applied before and after estimation.}\label{tab:tests}
\begin{tabularx}{\textwidth}{@{}>{\raggedright\arraybackslash}p{6.2cm} X@{}}
\toprule
Test & What it checks \\
\midrule
\multicolumn{2}{@{}l}{\emph{Pre-estimation (identification conditions)}}\\
Skewness, kurtosis, Jarque--Bera \citep{jarquebera1987} & Assumption~\ref{ass:nongauss}, on $\tx$ and on OLS residuals \\
Identification strength $D$, eq.~\eqref{eq:D} & Whether the higher-order information separates the sources \\
Breusch--Pagan / White \citep{breusch1979,white1980} & Heteroskedasticity, needed for the Lewbel competitor and for (L2) \\
\midrule
\multicolumn{2}{@{}l}{\emph{Relevance}}\\
Robust first-stage $F$ / Kleibergen--Paap \citep{montielolea2013,kleibergen2006} & Strength of $\xexo$; computed on $B$ only \\
Anderson--Rubin set \citep{andersonrubin1949} & Weak-instrument-robust inference on $\beta$ \\
\midrule
\multicolumn{2}{@{}l}{\emph{Exogeneity of $\xexo$ (the critical point)}}\\
HSIC / distance covariance of $\xexo$ and $\hat\varepsilon$ \citep{gretton2008,szekely2007} & Dependence beyond correlation (which is zero by construction) \\
Hansen $J$ with $f_1(\xexo),f_2(\xexo),\dots$ \citep{hansen1982} & Contaminated instruments imply different $\beta$s across functions \\
C test \citep{eichenbaum1988} & Validity of a subset of instruments \\
Agreement between strategies & Hausman-type contrast with Lewbel, SP, copula, Cumulant-ICA \\
Hausman, OLS vs IV \citep{hausman1978} & Is there endogeneity to correct? \\
Negative-control outcomes & Zero effect on an outcome $x$ cannot affect \\
$\HSIC(\xexo,r)$ & Independence of the split (necessary only) \\
\midrule
\multicolumn{2}{@{}l}{\emph{Sensitivity}}\\
\citet{conley2012}; \citet{nevo2012} & Range of $\hat\beta$ if $\Cov(\xexo,u)\ne0$ \\
\bottomrule
\end{tabularx}
\end{table}

\paragraph{Exogeneity cannot be tested directly.} Because $u$ is unobserved, every exogeneity test
uses the estimated residual $\hat\varepsilon = y - x\hat\beta - w'\hat\delta$ with $\hat\beta$
computed from $\xexo$ itself. The linear correlation between $\xexo$ and $\hat\varepsilon$ is
therefore zero by construction and only non-linear dependence is informative. Since $\hat\beta$ is
estimated, permutation or naive bootstrap calibration is only approximate; a bootstrap that
re-estimates $\beta$ in each draw is preferable.

\paragraph{Over-identification with functions of the instrument.} If $\xexo$ is contaminated, the
moment conditions $\E[f_j(\xexo)(y-x\beta-w'\delta)]=0$ generally imply different values of $\beta$
for different $f_j$, and Hansen's $J$ rejects. Two blind spots deserve emphasis, both confirmed by
the accompanying unit tests: under Gaussianity $\E[u\mid\xexo]$ is linear and all $f_j$ imply the
same wrong $\beta$; and if $H$ and $e_X$ are \emph{exchangeable} (identically distributed), then
$\E[H\mid H+e_X] = (H+e_X)/2$ is again linear and the test has no power against the fully
contaminated instrument $\tx$. The test has power when the conditional expectation of the
confounder given the instrument is non-linear.

\paragraph{Procedural safeguards.}
\begin{itemize}
  \item \emph{Circularity.} If the network minimised $\HSIC(\hat\varepsilon,\xexo)$, testing the same
  quantity on the same data tests what was optimised. Train on $A$, test on $B$.
  \item \emph{Gaussianity.} Under normality, independence reduces to zero correlation, which holds
  by construction, and every independence test passes vacuously. Always report the
  non-Gaussianity diagnostics alongside.
  \item \emph{Multiple testing.} When several exogeneity tests are reported, adjust with
  Bonferroni or \citet{holm1979}.
  \item \emph{Post-selection inference.} Hyper-parameters and architecture choices informed by test
  outcomes must be made on $A$.
\end{itemize}

\section{Related literature}\label{sec:lit}

Table~\ref{tab:lit} positions the proposal. Two strands are closest. \citet{dong2026sp} also split
$x$ without an external instrument, constructing $z = x - P_S x$ by projecting $x$ on a sieve
approximation of the confounding $\E[\varepsilon\mid x]$; identification there comes from functional
form (the confounding has no linear part in $x$), which fails under joint normality and does not
cover linear confounding. \citet{delbianco2026} tests the exclusion restriction of a single
instrument with fourth-order cumulants under a latent non-Gaussian factor, an over-complete
$2\times3$ ICA problem; the direct effect is identified only up to two points, which is the
counterpart of our Proposition~\ref{prop:twopoint}(b), and power requires substantial kurtosis and
$n$ in the low thousands.

\begin{table}[ht]
\centering
\small
\caption{Position in the literature.}\label{tab:lit}
\begin{tabularx}{\textwidth}{@{}>{\raggedright\arraybackslash}p{3.6cm} X X@{}}
\toprule
Work & What it does & Relation to this proposal \\
\midrule
\citet{dong2026sp} & Internal instrument $z = x - P_Sx$; identification by functional form; LASSO over a basis; non-parametric Hausman-type test & Splits $x$ without IV; different identification; no networks; fails under joint normality and linear confounding \\
\citet{delbianco2026} & Tests exclusion with one instrument via fourth-order cumulants; over-complete ICA; effect identified up to two points & Shows $e_X$ identifiable under latent confounding; needs $Z$, linear, power only with kurtosis and large $n$; provides a $D$ diagnostic \\
\citet{tramontano2025} & Causal effects in lvLiNGAM from higher-order cumulants with one proxy or instrument & Close neighbour of the identification layer \\
\citet{hoyer2008,salehkaleybar2020,wang2023,silva2017} & LiNGAM with latent variables; search for valid instruments under non-Gaussianity & Main risk to the novelty of the identification layer \\
\citet{saengkyongam2022} & HSIC-X: minimise HSIC between residuals and instruments & Source of the independence loss; still requires instruments \\
\citet{hyvarinen2019,khemakhem2020,zheng2022} & Identifiable non-linear ICA with auxiliary variables or sparsity & Conditions to make a non-linear split identifiable \\
\citet{lewbel1997,erickson2002} & Higher-moment instruments for mismeasured regressors & Ancestors of the cumulant layer \\
\citet{lewbel2012,parkgupta2012,ebbes2005,rigobon2003} & Internal instruments via heteroskedasticity, copulas, latent mixtures, regimes & Competitors and robustness checks \\
\citet{conley2012,nevo2012} & Inference with imperfect exogeneity & Sensitivity analysis \\
\citet{hartford2017,bennett2019} & Deep IV and deep GMM & Neural IV estimators that assume an external instrument \\
\bottomrule
\end{tabularx}
\end{table}

\paragraph{Apparent gap.} A non-linear split of the endogenous regressor, learned by adversarial or
HSIC independence, identified by non-Gaussianity, and used as an instrument or control function
with valid inference. The linear identification layer exists already; the defensible contribution
is the non-linear estimation layer with inference and Monte Carlo evidence. Referees will likely
describe the method as ``over-complete ICA by another name''; Proposition~\ref{prop:nosplit} and the
anchored split are our answer to what the network adds.

\section{Monte Carlo design}\label{sec:mc}

\subsection{Base DGP}

The base design is \eqref{eq:canonical} with $w\sim N(0,I_2)$, $\delta = (0.5,0.5)'$, $\beta=1$, unit
variances for all sources, and $g$ either linear or non-linear (softplus and sine terms). Sources are
drawn from Laplace (excess kurtosis $3$), centred exponential (skewness $2$), a symmetric bimodal
mixture (negative excess kurtosis), uniform, or Gaussian (negative control).

\subsection{Scenarios}

\begin{enumerate}[label=S\arabic*.]
  \item \emph{Non-Gaussian latent factor} (Delbianco's home turf): identification should work.
  Variants with linear and non-linear $g$.
  \item \emph{Non-linear confounding without a linear part} (SP's home turf):
  $x = g(w) + v$, $u = \lambda(v^2 - \E v^2) + e_Y$ with skewed $v$, so OLS is biased while the
  confounding is a quadratic function of $x$. SP should win; the proposal should at least not
  lose badly, and its diagnostics should flag the misspecification.
  \item \emph{Gaussian sources} (negative control): no non-Gaussian method identifies $\beta$; the
  procedure must report weak identification rather than a confident $\hat\beta$.
  \item \emph{Violated independence}: correlation $\rho$ between $H$ and $e_X$, or a direct channel
  $u \mathrel{+}= \delta_{\mathrm{excl}} e_X$. Measures bias as the hypothesis breaks.
  \item \emph{Placebo and size}: $\lambda=0$ (exogenous $x$) and pure-noise candidate instruments.
  The procedure must not invent exogenous variation nor ``pass'' tests with garbage.
  \item[Aux.] \emph{Auxiliary heteroskedasticity}: $\Var(e_X\mid w_1)$ varies with $w_1$, which
  enables labelling (L2) and gives the Lewbel competitor its best case.
\end{enumerate}

\subsection{Grid, estimators and metrics}

The full grid uses $n\in\{500,1000,2000,5000,10000\}$, $\lambda\in\{0.2,0.5,1\}$ and 500--1000
replications per cell. Estimators: OLS; 2SLS with the planted exogenous variation $e_X$ as an
oracle instrument (performance ceiling); SP (sieve + LASSO); Lewbel; Gaussian-copula control
function; Cumulant-ICA with labelling (L1) and (L2); and the neural split (IV and control-function
versions). Delbianco's test is run only in an auxiliary design with a planted instrument. Metrics:
bias, median bias and RMSE of $\hat\beta$; coverage of nominal 95\% intervals (Wald, bootstrap or
Anderson--Rubin); size and power of each test of Section~\ref{sec:tests} against planted violations;
and the share of replications in which the method declares that $\beta$ is not identified.

\subsection{Pilot results}\label{sec:pilot}

Tables~\ref{tab:pilot-fast} and~\ref{tab:pilot-nn} report a small pilot run
($\lambda = 0.5$, labelling (L1) with $\lambda>0$ known) produced by
\texttt{scripts/run\_montecarlo.py}; the full grid is left for the next stage.

\textit{[pilot_fast pending]}


\textit{[pilot_nn pending]}


\textit{[pilot_discussion pending]}


\section{Risks and limitations}\label{sec:risks}

\begin{itemize}
  \item Assumption~\ref{ass:indep} is equivalent to a disguised exclusion restriction and is not
  testable. The paper must say so plainly, and results should always be accompanied by sensitivity
  analysis and by at least one alternative identification strategy.
  \item Identification through higher-order moments is weak: it requires non-trivial kurtosis (or
  skewness) and large samples; $D$ makes this visible but does not cure it.
  \item The label ambiguity of Proposition~\ref{prop:twopoint}(b) is intrinsic. Assumption~\ref{ass:label}
  is an additional, substantive input.
  \item Over-complete non-linear ICA is less identifiable than its linear counterpart; without a
  class restriction the split is arbitrary. Proposition~\ref{prop:nosplit} shows that even in the
  linear case the exogenous source is not a function of the data.
  \item Novelty risk: the identification layer already exists in the latent-variable LiNGAM
  literature; the contribution must lie in non-linear estimation, inference and simulation evidence.
\end{itemize}

\section{Next steps}

\begin{enumerate}
  \item Run scenarios S1 and S3 first: if the method cannot separate the identifiable case from the
  Gaussian one, nothing else matters. The pilot shows that the cumulant layer and the diagnostic $D$
  pass this test, while the unanchored network does not improve on OLS.
  \item Implement and evaluate the anchored split of Section~\ref{sec:anchored}.
  \item Replace permutation calibration of the residual-independence tests with a bootstrap that
  re-estimates $\beta$.
  \item Read \citet{hoyer2008}, \citet{silva2017} and \citet{tramontano2025} closely to settle the
  novelty question.
\end{enumerate}

\appendix
\section{Implementation and tests}\label{app:code}

The accompanying Python package \texttt{ivsplit} implements every component of the paper:

\begin{center}
\small
\begin{tabularx}{\textwidth}{@{}l X@{}}
\toprule
Module & Content \\
\midrule
\texttt{dgp.py} & Source distributions, DGP \eqref{eq:canonical}, scenarios S1--S5 and Aux \\
\texttt{estimators.py} & OLS, 2SLS (robust), oracle, Lewbel, copula CF, SP, joint cumulants, Prony roots, labelling, Cumulant-ICA \\
\texttt{independence.py} & HSIC, distance covariance, permutation tests \\
\texttt{split.py} & Neural split, loss \eqref{eq:loss}, sample-split estimation and inference \\
\texttt{diagnostics.py} & The battery of Table~\ref{tab:tests}, Conley et al.\ bounds, Holm/Bonferroni \\
\texttt{montecarlo.py} & Grid runner and summary metrics \\
\bottomrule
\end{tabularx}
\end{center}

The test suite (\texttt{pytest}) checks, among other things: that sources are standardised and have
the intended skewness and kurtosis; that the oracle 2SLS is consistent while OLS converges to
$\beta + \lambda\sigma_H^2/(\sigma_H^2+\sigma_X^2)$; that the Prony roots recover $\{\beta,\beta+\lambda\}$
and flipping the sign restriction returns $\beta+\lambda$; that $D$ is large under non-Gaussian
confounding and small under Gaussian sources or no confounding; that SP wins in S2 and fails under
linear confounding; that the over-identification test with powers of the instrument detects
contamination when $\E[H\mid\tx]$ is non-linear and is blind when the sources are Gaussian or
exchangeable; that HSIC and distance covariance detect non-linear dependence and have correct size;
and that the neural split respects the relevance band, uses disjoint halves, and reports failed
identification under Gaussian sources.

\begingroup\sloppy
\bibliographystyle{plainnat}
\bibliography{references}
\endgroup

\end{document}
